\chapter{Εισαγωγή}
\label{ch:introduction}

\section{Το πλαίσιο: Ολοκληρωμένα Δορυφορικά-Επίγεια Δίκτυα στην εποχή του 6G}
\label{sec:intro-context}

Η ζήτηση για ασύρματη συνδεσιμότητα υψηλού ρυθμού και καθολικής κάλυψης
αυξάνεται συνεχώς, καθώς εφαρμογές έντασης δεδομένων, όπως η ροή βίντεο,
οι υπηρεσίες υπολογιστικού νέφους και το Διαδίκτυο των Πραγμάτων (IoT),
γίνονται ολοένα πιο απαιτητικές. Τα συμβατικά επίγεια δίκτυα
(Terrestrial Networks - TN), παρά τη συνεχή εξέλιξή τους από τη 2η έως
την 5η γενιά κινητών επικοινωνιών, παραμένουν περιορισμένα από το κόστος
και την πολυπλοκότητα ανάπτυξης υποδομής σε αγροτικές, απομακρυσμένες ή
δύσβατες περιοχές, καθώς και σε θαλάσσιες και εναέριες περιοχές όπου η
εγκατάσταση σταθμών βάσης είναι απλά ανέφικτη. Η επικοινωνία μέσω
δορυφόρων (Satellite Communication - SatCom) προσφέρει, αντίθετα, σχεδόν
καθολική γεωγραφική κάλυψη, αλλά ιστορικά υστερούσε σε ρυθμό μετάδοσης,
καθυστέρηση και κόστος ανάπτυξης και συντήρησης σε σχέση με τα επίγεια
δίκτυα.

Η μείωση του κόστους εκτόξευσης και η εμπορική ανάπτυξη πυκνών
αστερισμών δορυφόρων χαμηλής τροχιάς (Low Earth Orbit - LEO) τα
τελευταία χρόνια κατέστησαν εφικτή την ενσωμάτωση μη-επίγειων δικτύων
(Non-Terrestrial Networks - NTN) στα πρότυπα κινητών επικοινωνιών. Το
3GPP συμπεριέλαβε επίσημα την υποστήριξη NTN στο Release~17 του 5G New
Radio (NR)~\cite{tr38821}, ανοίγοντας τον δρόμο για δίκτυα που
συνδυάζουν επίγειους σταθμούς βάσης και δορυφόρους σε ένα ενιαίο,
τρισδιάστατο δίκτυο πρόσβασης. Τέτοια δίκτυα αναφέρονται στη διεθνή
βιβλιογραφία ως Ολοκληρωμένα Δορυφορικά-Επίγεια Δίκτυα
(Satellite-Terrestrial Integrated Networks - STINs) και αποτελούν έναν
από τους βασικούς σχεδιαστικούς στόχους του 6G~\cite{li2025advancing,
noh2024tnntn}.

Η απλή συνύπαρξη TN και NTN δεν αρκεί ωστόσο για να αξιοποιηθούν πλήρως
τα συμπληρωματικά τους πλεονεκτήματα. Όταν ένας χρήστης μπορεί να
συνδεθεί είτε σε επίγειο σταθμό βάσης είτε σε δορυφόρο, αλλά επιλέγει
πάντα μόνο τον έναν από τους δύο (single-connectivity), το δίκτυο δεν
εκμεταλλεύεται τη δυνατότητα ταυτόχρονης ή δυναμικά εναλλασσόμενης
χρήσης και των δύο ζεύξεων. Η \emph{πολυσυνδεσιμότητα} (Multi-Connectivity
- MC) -- η ταυτόχρονη ή συντονισμένη σύνδεση ενός χρήστη σε περισσότερους
από έναν κόμβους δικτύου -- έχει προταθεί ως μηχανισμός βελτίωσης της
αξιοπιστίας, της κάλυψης, της διαχείρισης φορτίου και της ενεργειακής
απόδοσης σε STINs, και παραμένει ενεργό αντικείμενο έρευνας τόσο στη
βιβλιογραφία όσο και στο ίδιο το 3GPP~\cite{majamaa2024toward,
pupiales2021multiconnectivity, shang2024multiconnectivity}.

Η ενσωμάτωση πολυσυνδεσιμότητας σε ένα σύστημα όπου ο ένας από τους δύο
υποψήφιους κόμβους (ο δορυφόρος LEO) κινείται με ταχύτητα ίχνους εδάφους
της τάξης μερικών km/s φέρνει, ωστόσο, προκλήσεις άγνωστες στα καθαρά
επίγεια δίκτυα: το κανάλι προς τον δορυφόρο αλλάζει γρήγορα με τον χρόνο,
οι αποφάσεις επιλογής ή εναλλαγής κόμβου (handover) πρέπει να λαμβάνονται
συχνά και υπό αβεβαιότητα, και η ισορροπία μεταξύ ρυθμού μετάδοσης και
κατανάλωσης ισχύος διαφέρει ριζικά ανάμεσα σε επίγεια και δορυφορική
ζεύξη. Το σημείο εκκίνησης της παρούσας εργασίας είναι ακριβώς η
κατανόηση αυτής της δυναμικής μέσω προσομοίωσης, πριν από οποιαδήποτε
προσπάθεια βελτιστοποίησης ή αυτοματοποίησης με τεχνικές μηχανικής
μάθησης.

\section{Αντικείμενο διπλωματικής}
\label{sec:intro-objective}

Η παρούσα διπλωματική εργασία αναπτύσσει και αξιολογεί ένα πλαίσιο
προσομοίωσης σε MATLAB για ένα ελάχιστο αλλά ρεαλιστικό σενάριο STIN: δύο
επίγειοι σταθμοί βάσης 5G NR και ένας δορυφόρος LEO εξυπηρετούν από κοινού
ένα σύνολο χρηστών, καθένας εκ των οποίων συνδέεται σε κάθε υποψήφιο κόμβο
(επίγειο σταθμό βάσης, δορυφόρο) του οποίου το SNR ξεπερνά ένα ελάχιστο
χρησιμοποιήσιμο κατώφλι -- ταυτόχρονα και στους δύο όταν και οι δύο το
ξεπερνούν (πολυσυνδεσιμότητα, greedy per-user απόφαση). Συγκεκριμένα, η
εργασία:

\begin{enumerate}
  \item Μοντελοποιεί το επίγειο κανάλι σύμφωνα με το πρότυπο 3GPP
    TR~38.901~\cite{tr38901} (σενάρια UMa/UMi, στοχαστική επιλογή
    LOS/NLOS ανά ζεύξη βάσει της απόστασης, log-normal shadow fading), και
    το δορυφορικό κανάλι με απώλειες ελεύθερου χώρου (FSPL) υπό μάσκα
    ελάχιστης γωνίας ανύψωσης.
  \item Υπολογίζει, ανά χρήστη, δείκτες απόδοσης (Key Performance
    Indicators - KPIs): ρυθμό μετάδοσης (χωρητικότητα Shannon μετά από
    ισομερή κατανομή εύρους ζώνης), και ενέργεια ανά bit, βάσει γραμμικών
    μοντέλων κατανάλωσης ισχύος (μοντέλο EARTH~\cite{auer2011earth} για
    τον σταθμό βάσης, μοντέλο απόδοσης ενισχυτή ισχύος για τον δορυφόρο).
  \item Επεκτείνει το εγγενώς στατικό σενάριο αναφοράς προς δύο
    κατευθύνσεις που το πλησιάζουν σε ρεαλιστική λειτουργία: πολλαπλές
    εκτελέσεις της ίδιας τοπολογίας με νέα τυχαία πραγματοποίηση καναλιού
    κάθε φορά, ώστε να χαρακτηριστεί η διακύμανση των KPI, και προσομοίωση
    διαδοχικών μεταδόσεων κατά τη διάρκεια μίας διέλευσης του δορυφόρου LEO,
    με το υποδορυφορικό σημείο να κινείται με πραγματική ταχύτητα ίχνους
    εδάφους, υπολογισμένη από την τροχιακή περίοδο Kepler.
  \item Παράγει, μέσω ενός Monte Carlo driver πάνω σε τυχαιοποιημένες
    τοπολογίες, ένα ετικετοποιημένο σύνολο δεδομένων, και εκπαιδεύει τρεις
    παραλλαγές δύο μοντέλων μηχανικής μάθησης (Λογιστική Παλινδρόμηση,
    Random Forest) που προβλέπουν την κατάσταση σύνδεσης
    (\texttt{ServingType}: επίγειος/δορυφορικός/ταυτόχρονα-και-τα-δύο)
    από τα χαρακτηριστικά των υποψήφιων ζεύξεων, με προοδευτικά λιγότερη,
    πιο ρεαλιστική πληροφορία SNR διαθέσιμη στο μοντέλο σε κάθε παραλλαγή
    (Ενότητα~\ref{sec:results-ml-realistic}): από το ακριβές, στιγμιαίο
    SNR (μη ρεαλιστικό, χρήσιμο μόνο για έλεγχο ορθότητας του pipeline),
    μέσω ενός θορυβώδους SNR που προσομοιώνει μια πραγματική, ατελή
    μέτρηση (το σενάριο που η εργασία θεωρεί αντιπροσωπευτικό μιας
    πραγματικής ανάπτυξης), έως την πλήρη απουσία SNR (μόνο γεωμετρία).
\end{enumerate}

Ο απώτερος στόχος, πέραν της παρούσας εργασίας, είναι η ανάπτυξη ενός
μοντέλου μηχανικής μάθησης που θα υποκαθιστά τον απλό κανόνα σύγκρισης
SNR με μια πιο ρεαλιστική απόφαση επιλογής ή εναλλαγής κόμβου, ανθεκτική
σε μερική ή θορυβώδη πληροφορία. Η παρούσα εργασία θέτει τα θεμέλια
αυτής της προσπάθειας, και κάνει ήδη ένα πρώτο βήμα προς αυτήν: τον
προσομοιωτή που παράγει τα δεδομένα εκπαίδευσης, την ποσοτικοποίηση της
συμπεριφοράς του σε στατικά και χρονικά μεταβαλλόμενα σενάρια, και τη
σύγκριση μοντέλων μηχανικής μάθησης υπό ρεαλιστικά διαθέσιμη πληροφορία
αντί για ιδανική (βλ.~Κεφάλαιο~\ref{ch:results} και
Κεφάλαιο~\ref{ch:discussion}).

\subsection{Ερευνητικά ερωτήματα}
\label{subsec:intro-questions}

Η εργασία διαρθρώνεται γύρω από τα εξής ερωτήματα:

\begin{itemize}
  \item[\textbf{Ε1.}] Ποια είναι η αναμενόμενη απόδοση (throughput, ενέργεια
    ανά bit, SNR) ενός χρήστη όταν εξυπηρετείται από επίγειο σταθμό βάσης
    έναντι δορυφόρου LEO, σε ένα ρεαλιστικό μοντέλο καναλιού 3GPP;
  \item[\textbf{Ε2.}] Πώς μεταβάλλονται αυτοί οι δείκτες, και η ίδια η
    απόφαση επιλογής κόμβου, καθώς ο δορυφόρος LEO διασχίζει τον ουρανό
    πάνω από την περιοχή κάλυψης κατά τη διάρκεια μίας διέλευσης;
  \item[\textbf{Ε3.}] Ποιοι περιορισμοί του τρέχοντος (greedy, στιγμιαίας
    σύγκρισης SNR) κανόνα επιλογής κόμβου αναδεικνύονται όταν η
    προσομοίωση επεκτείνεται από ένα μεμονωμένο στατικό σενάριο σε
    διαδοχικές μεταδόσεις με κινούμενο δορυφόρο, και τι θα χρειαστεί
    ώστε ένα μελλοντικό μοντέλο μηχανικής μάθησης να τους αντιμετωπίσει;
\end{itemize}

\section{Συνεισφορά}
\label{sec:intro-contribution}

Η συνεισφορά της εργασίας συνοψίζεται ως εξής:

\begin{itemize}
  \item Ένα πλήρες, αναπαραγώγιμο (fixed RNG seed) πλαίσιο προσομοίωσης
    link-budget για σενάρια STIN με δύο επίγειους σταθμούς βάσης και έναν
    δορυφόρο LEO, βασισμένο σε πρότυπα 3GPP και στο 5G Toolbox της
    MATLAB.
  \item Στατιστική ανάλυση τριών KPI (throughput, ενέργεια ανά bit, SNR)
    πάνω από πολλαπλές τυχαίες πραγματοποιήσεις μιας σταθερής τοπολογίας
    αναφοράς.
  \item Μια επέκταση της προσομοίωσης σε \emph{χρονική} διάσταση, απούσα
    από την αρχική, στατική της έκδοση, με ρεαλιστική κίνηση του
    δορυφόρου κατά τη διάρκεια μίας διέλευσης, και καταγραφή handovers
    και συμπεριφοράς KPI με τον χρόνο.
  \item Ένα χωρικά συσχετισμένο μοντέλο shadow fading κατά Gudmundson
    (Ενότητα~\ref{subsec:meth-correlated-sf}), που διόρθωσε ένα πρόβλημα
    το οποίο ανέδειξε ακριβώς αυτή η χρονική επέκταση: μη ρεαλιστικά
    συχνά handovers, οφειλόμενα σε ανεξάρτητο δειγματισμό σκίασης ανά
    χρονικό βήμα.
  \item Τρεις παραλλαγές μοντέλων μηχανικής μάθησης, με προοδευτικά
    λιγότερη πληροφορία SNR: μία που επαληθεύει την ορθότητα του
    συνολικού pipeline (MATLAB $\rightarrow$ CSV $\rightarrow$ Python) με
    ακριβές SNR (μη ρεαλιστική ως στόχος, χρήσιμη μόνο ως έλεγχος), μία
    με θορυβώδες SNR που αναπαριστά μια ρεαλιστική, ατελή μέτρηση --
    αυτή θεωρείται η αντιπροσωπευτική επίδοση της εργασίας -- και μία
    χωρίς καθόλου SNR, μόνο γεωμετρία, ως κατώτατο όριο σύγκρισης.
  \item Μια ρητή κατάσταση outage στον κανόνα επιλογής κόμβου
    (εξίσωση~\eqref{eq:node-selection}), που διόρθωσε ένα δεύτερο
    πρόβλημα ανιχνευμένο από την ίδια χρονική επέκταση: χρήστες χωρίς
    κανέναν πρακτικά χρησιμοποιήσιμο υποψήφιο ανατίθονταν στον "λιγότερο
    κακό" κόμβο αντί να καταγράφονται ως εκτός κάλυψης.
  \item Επέκταση του κανόνα επιλογής κόμβου σε πραγματική
    πολυσυνδεσιμότητα (εξίσωση~\eqref{eq:node-selection}): ένας χρήστης
    εξυπηρετείται πλέον ταυτόχρονα από BS \emph{και} δορυφόρο όταν και οι
    δύο ξεπερνούν το ελάχιστο χρησιμοποιήσιμο SNR (το τετρα-καταστασιακό
    μοντέλο SS-SBS της Ενότητας~\ref{subsec:bg-architectures}), με
    αθροιστική χωρητικότητα/ενέργεια στις δύο ενεργές ζεύξεις. Η αλλαγή
    αυτή ανέδειξε με τη σειρά της ένα εύρημα για το ίδιο το σενάριο
    προσομοίωσης (Ενότητα~\ref{sec:results-ml-realistic}): με τη
    γεωμετρία δορυφόρου/κάλυψης του σεναρίου αναφοράς, η ταυτόχρονη
    σύνδεση και στους δύο κόμβους αποδεικνύεται ο κανόνας παρά η
    εξαίρεση, γεγονός που άλλαξε αισθητά και τη συμπεριφορά των τριών
    πειραμάτων μηχανικής μάθησης παραπάνω.
  \item Μια μηχανή υστέρησης (hysteresis) και time-to-trigger στην
    ενεργοποίηση/απενεργοποίηση κάθε ζεύξης (εξίσωση~\eqref{eq:hysteresis}),
    κατά το πρότυπο Event~A3 του 3GPP TS~38.331. Επειδή το ίδιο το
    σενάριο αναφοράς αποδείχθηκε "ήσυχο" (καθόλου ταλάντωση κοντά στο
    κατώφλι σε ακίνητους χρήστες), ο μηχανισμός επικυρώθηκε ξεχωριστά σε
    ένα ειδικά κατασκευασμένο σενάριο πίεσης: $67\%$ μείωση τεχνητών
    εναλλαγών κατάστασης ($12\rightarrow4$ σε $80$ βήματα), με εμφανή
    καταστολή ping-pong. Η joint/fairness-aware πολυσυνδεσιμότητα ορίζει
    πλέον την επόμενη προτεραιότητα της εργασίας.
\end{itemize}

\section{Δομή της διπλωματικής}
\label{sec:intro-structure}

Το Κεφάλαιο~\ref{ch:background} παρουσιάζει το θεωρητικό υπόβαθρο και τη
σχετική βιβλιογραφία: αρχιτεκτονικές πολυσυνδεσιμότητας σε STINs, τις
προκλήσεις σχεδίασής τους, την πορεία τυποποίησης TN-NTN στο 3GPP, και τα
μοντέλα διάδοσης καναλιού στα οποία βασίζεται η προσομοίωση. Το
Κεφάλαιο~\ref{ch:methodology} περιγράφει αναλυτικά τη μεθοδολογία: το
μοντέλο επίγειου και δορυφορικού καναλιού, τον κανόνα επιλογής κόμβου, τα
μοντέλα χωρητικότητας και ενέργειας, και το μοντέλο κίνησης του δορυφόρου
που αναπτύχθηκε για τη χρονική προσομοίωση. Το
Κεφάλαιο~\ref{ch:implementation} περιγράφει την υλοποίηση σε λογισμικό
(δομή κώδικα, ροή δεδομένων, αναπαραγωγιμότητα). Το
Κεφάλαιο~\ref{ch:results} παρουσιάζει τα αποτελέσματα: το σενάριο
αναφοράς, τη στατιστική ανάλυση KPI, την προσομοίωση διέλευσης, και την
αξιολόγηση των τριών παραλλαγών μηχανικής μάθησης (ακριβές, θορυβώδες,
χωρίς SNR). Το
Κεφάλαιο~\ref{ch:discussion} σχολιάζει τα ευρήματα, καταγράφει τους
περιορισμούς του τρέχοντος μοντέλου, και σχεδιάζει τη μελλοντική εργασία.
Το Κεφάλαιο~\ref{ch:conclusions} συνοψίζει τα συμπεράσματα. Το
Παράρτημα~Ι περιλαμβάνει επιλεγμένα τμήματα κώδικα που υλοποιούν τον
πυρήνα της προσομοίωσης.
